\subsubsection{Cutout/Coarse Dropout} \label{Cutout} 
Cutout is an image augmentation technique that involves randomly selecting a rectangular region of an image and setting the pixel values in that region to zero or a constant value. This creates a "hole" in the image, which can be used to encourage the model to focus on the remaining information in the image and improve its robustness to occlusion or missing data. The size, shape, and number of cutout regions can be varied to create a diverse set of training images and prevent overfitting.

Cutout augmentation can potentially be non label-preserving in aerial images because it involves removing parts of the image, which can also remove relevant objects or features that the model needs to learn, e.g a bird in the center of the image. This can cause inconsistencies between the original labels of the image and the labels assigned to the modified image, especially if the removed region contains the object of interest. For examples in Figure\ref{fig:Cutout}, one common goldeneye got cut out, if there were just one bird in the picture, it would be considered non-label preserving.

\begin{figure}[!h]
	 \centerline{\includegraphics[width=0.8\textwidth]{src/augmentations/Cutout.png}}
	 {\caption{Examples of Cutout with 8 boxes, each size 1/10 height of the image} \label{fig:Cutout}}
\end{figure}

\subsubsection{GridDropout}

GridDropout is an image augmentation technique that is based on the Cutout method \ref{Cutout} by the same author \cite{chen2020gridmask}. It involves dividing the image into a grid of equal-sized cells and randomly dropping out entire cells instead of rectangular patches. This creates a checkerboard pattern of dropped cells that encourages the model to learn features that are invariant to the exact location of objects in the image. The advantage of using GridDropout over Cutout is that it allows for larger and more complex holes to be created while also being less likely to remove entire objects from the image. In the case of our aerial image dataset, the objects of interest are often small in size, and therefore applying cutout augmentation may not be label-preserving due to the risk of removing important object features, as illustrated in Figure\ref{fig:GridDropout}.

\begin{figure}[!h]
	 \centerline{\includegraphics[width=0.8\textwidth]{src/augmentations/GridDropout.png}}
	 {\caption{Examples of GridDropout with 100 boxes} \label{fig:GridDropout}}
\end{figure}